\begin{lemma}[The Type~I infimum is attained, at a pair of at most half the total length]
  \label{typeI-cut-minimal-pair}
  Let $E$ be a metric space and $\gamma \in \Theta(E)$ (\noderef{Curve}) a closed curve
  (\noderef{IsClosedCurve}); write $L := \length(\gamma)$. Let $\delta,\epsilon \in \mathbb{R}$
  with $\delta > 0$, and for $a,b \in \mathbb{R}$ set
  \[
    \ell(a,b) := \begin{cases} b-a & a \le b \\ (L-a)+b & a > b \end{cases}
    .
  \]
  Call a pair $(a,b) \in [0,L]^2$ \emph{admissible} if
  \[
    \delta \le d_\gamma(a,b)
    \qquad\text{and}\qquad
    d(\gamma(a),\gamma(b)) \le \epsilon\, d_\gamma(a,b)
  \]
  (\noderef{dGamma}), and suppose some pair $(a_0,b_0) \in [0,L]^2$ is admissible. Then there is an
  admissible pair $(t,t') \in [0,L]^2$ with
  \[
    \delta \le \ell(t,t') \le \tfrac{1}{2}L
  \]
  which minimizes $\ell$ among admissible pairs: $\ell(t,t') \le \ell(a,b)$ for every admissible
  $(a,b) \in [0,L]^2$.

  \paragraph{Relation to the paper.} This is the attainment claim of paper.tex 857-864 together
  with the two bounds $\beta(\gamma) \ge \delta$ (paper.tex 859) and $\beta(\gamma) \le
  \tfrac12\length(\gamma)$ (paper.tex 860-861), for the infimum \eqref{betagammainf}: the set of
  admissible pairs is exactly the index set of that infimum, and $\ell(a,b)$ is exactly the
  quantity $\length(\gamma|_{[a,b]})$ it minimizes (that identification is
  \Cref{curve-arc-eval}, and is not needed here, where $\ell$ is taken as the explicit formula
  above). Non-emptiness of the index set is a \emph{hypothesis} here, supplied by the witness
  $(a_0,b_0)$; in the paper it comes from the failure of large-scale invertibility, which is the
  consumer's business, not this lemma's. No claim of definiteness is made about $(t,t')$: the
  paper's choice of the lexicographically-smallest attaining pair carries no mathematical content
  beyond attainment.

  \paragraph{No order relation between $t$ and $t'$.} The minimizing pair may well have $t' < t$,
  i.e.\ the counter-clockwise arc from $t$ to $t'$ may be the wrap-around one; nothing in the
  statement singles out either order, and the bound $\ell(t,t') \le \tfrac12 L$ is precisely what
  survives in place of an order assumption.
\end{lemma}
\begin{proof}
  Write $S \subseteq \mathbb{R}^2$ for the set of admissible pairs and
  \[
    D(a,b) := \min(|a-b|,\, L-|a-b|)
    .
  \]
  Since $\gamma$ is closed, $d_\gamma(a,b) = D(a,b)$ for all $a,b$ (\noderef{dGamma}: the closed
  branch of its defining case split applies), so $S = \{(a,b) \in [0,L]^2 : \delta \le D(a,b)
  \text{ and } d(\gamma(a),\gamma(b)) \le \epsilon D(a,b)\}$, and we work with $D$ from now on.

  \paragraph{Four elementary facts about $\ell$ and $D$.}

  \emph{(F1) $D(a,b) \le \ell(a,b)$ for all $a,b \in \mathbb{R}$.} If $a \le b$ then $|a-b| = b-a =
  \ell(a,b)$, and $D(a,b) \le |a-b|$ since $D$ is a minimum of two terms the first of which is
  $|a-b|$. If $a > b$ then $|a-b| = a-b$, so $L - |a-b| = (L-a)+b = \ell(a,b)$, and $D(a,b) \le
  L-|a-b|$.

  \emph{(F2) If $a \ne b$ then $\ell(a,b) + \ell(b,a) = L$.} If $a < b$ then $\ell(a,b) = b-a$ and
  $\ell(b,a) = (L-b)+a$, summing to $L$. If $a > b$ then $\ell(a,b) = (L-a)+b$ and $\ell(b,a) =
  a-b$, again summing to $L$.

  \emph{(F3) If $a \ne b$ and $\ell(a,b) \le \ell(b,a)$ then $D(a,b) = \ell(a,b)$.} If $a < b$ then
  $\ell(a,b) = b-a = |a-b|$ and $\ell(b,a) = (L-b)+a = L-|a-b|$, so the hypothesis reads $|a-b| \le
  L-|a-b|$ and $D(a,b) = \min(|a-b|,L-|a-b|) = |a-b| = \ell(a,b)$. If $a > b$ then $\ell(a,b) =
  (L-a)+b = L-|a-b|$ and $\ell(b,a) = a-b = |a-b|$, so the hypothesis reads $L-|a-b| \le |a-b|$ and
  $D(a,b) = L-|a-b| = \ell(a,b)$.

  \emph{(F4) $S$ is symmetric: $(a,b) \in S$ iff $(b,a) \in S$.} $|a-b| = |b-a|$, so $D(a,b) =
  D(b,a)$; and $d(\gamma(a),\gamma(b)) = d(\gamma(b),\gamma(a))$ by symmetry of the metric of $E$.
  Both defining conditions of $S$, and membership in $[0,L]^2$, are therefore symmetric.

  \paragraph{$S$ is compact and non-empty.} $S$ is the intersection of the square $[0,L]^2$ --
  compact, being a product of two compact intervals -- with the two sets $\{(a,b) : \delta \le
  D(a,b)\}$ and $\{(a,b) : d(\gamma(a),\gamma(b)) \le \epsilon D(a,b)\}$. The map $(a,b) \mapsto
  D(a,b)$ is continuous on $\mathbb{R}^2$, being built from the continuous operations
  $(a,b)\mapsto a-b$, $|\cdot|$, $x \mapsto L-x$ and $\min$. The underlying map of $\gamma$ is
  $1$-Lipschitz on all of $\mathbb{R}$ (\noderef{CurveLipschitz}), hence continuous, so $(a,b)
  \mapsto d(\gamma(a),\gamma(b))$ is continuous as a composition of continuous maps; and so is
  $(a,b) \mapsto \epsilon D(a,b)$. Each of the two sets is therefore closed, being of the form
  $\{f \le g\}$ for continuous $f,g$. Hence $S$, a compact set intersected with two closed sets, is
  compact. It is non-empty because $(a_0,b_0) \in S$ by hypothesis.

  \paragraph{Minimizing $D$ (not $\ell$) over $S$.} $D$ is continuous and $S$ is compact and
  non-empty, so by the extreme value theorem $D$ attains a minimum on $S$: fix $q = (u,v) \in S$
  with
  \[
    D(u,v) \le D(a,b) \qquad \text{for every } (a,b) \in S
    .
  \]
  (This is the step at which the present proof departs from the paper's, which minimizes
  $\ell$ directly. The paper's assertion at paper.tex 858 that $\length(\gamma|_{[s,s']})$ is a
  continuous function of $(s,s')$ is false -- $\ell$ jumps by $L$ across the diagonal $s=s'$ --
  so $\ell$ is not available to the extreme value theorem on the square. $D$ \emph{is} continuous
  everywhere, and by (F1) and (F3) below its minimum over $S$ is also the minimum of $\ell$ over
  $S$, once the minimizing pair is put in the right order.)

  Note $u \ne v$: indeed $\delta \le D(u,v) \le |u-v|$ by definition of $D$ as a minimum, and
  $\delta > 0$, so $|u-v| > 0$.

  \paragraph{Ordering the minimizing pair.} By (F4), $(v,u) \in S$ as well, and $D(v,u) = D(u,v)$.
  Define
  \[
    (t,t') := \begin{cases} (u,v) & \ell(u,v) \le \ell(v,u) \\ (v,u) & \text{otherwise.}\end{cases}
  \]
  In either case $(t,t') \in S$ (so $t,t' \in [0,L]$ and $(t,t')$ is admissible), $t \ne t'$,
  $D(t,t') = D(u,v)$, and
  \[
    \ell(t,t') \le \ell(t',t)
    ,
  \]
  the displayed inequality being the defining condition in the first case and, in the second case,
  the negation $\ell(v,u) < \ell(u,v)$ of that condition, read with the roles of the two arguments
  exchanged.

  \paragraph{The three conclusions.} By (F3) applied to $(t,t')$ -- legitimate since $t \ne t'$ and
  $\ell(t,t') \le \ell(t',t)$ -- we have
  \[
    D(t,t') = \ell(t,t')
    .
  \]
  Since $(t,t') \in S$, $\delta \le D(t,t') = \ell(t,t')$, which is the lower bound. By (F2),
  $\ell(t,t') + \ell(t',t) = L$, and $\ell(t,t') \le \ell(t',t)$, so $2\ell(t,t') \le L$, i.e.\
  $\ell(t,t') \le \tfrac12 L$, which is the upper bound. Finally, let $(a,b) \in [0,L]^2$ be
  admissible, i.e.\ $(a,b) \in S$. Then
  \[
    \ell(t,t') = D(t,t') = D(u,v) \le D(a,b) \le \ell(a,b)
    ,
  \]
  the middle inequality by minimality of $D$ at $q$ and the last by (F1). This is the asserted
  minimality of $\ell(t,t')$, and completes the proof.
\end{proof}
